\section{Round-sixteen positive closure: dominated information topology and stable belief dynamics}
\label{sec:r16-c1}

Weak convergence of joint laws does not make conditional probabilities
continuous.  The information state therefore retains the unnormalized joint
state--observation kernel in a topology which controls its density, rather than
reconstructing a posterior from a weak joint law.  Regular dominated charts,
singular charts, and zero-evidence directions are separately typed.

\subsection{The information state}

Let \(E\) be the microscopic or kinetic state and \(Y\) the observation space.
On a regular observation chart fix a sigma-finite reference measure \(\lambda\)
and write the one-step joint kernel as
\[
 K(\rho,dx,dy)=k_\rho(x,y)\,m(dx)\lambda(dy).
\]
For every Lyapunov sublevel \(\mathcal C_L\), assume the model-derived bounds
\[
 0<c_L\le g_\rho(y):=\int k_\rho(x,y)m(dx)\le C_L,
 \qquad
 \|k_\rho-k_{\tilde\rho}\|_{L^1_\lambda(BL_W^*)}
 \le C_Ld_{\mathcal I}(\rho,\tilde\rho).
\]
Here the mixed norm integrates, in \(y\), the weighted bounded-Lipschitz dual
norm of the unnormalized state measure.  The information metric is
\[
 d_{\mathcal I}(\rho,\tilde\rho)
 =d_{BL,W}(\rho,\tilde\rho)
 +\|k_\rho-k_{\tilde\rho}\|_{L^1_\lambda(BL_W^*)}.
\]
The posterior is the positive measure
\[
 B(\rho,y)(dx)={k_\rho(x,y)m(dx)\over g_\rho(y)}.
\]
At a singular observation stratum we retain its positive Rokhlin kernel as a
separate state component with its own dominating measure.  The zero-evidence
fiber stores the projective unnormalized kernel but is never used to divide by
zero.

\begin{theorem}[Model-derived dominated observation chart]
\label{thm:r16-c1-chart}
For every compact regular family of path/contact observation maps and compact
control set, the B1 exact coefficient and B2 inserted-source theorem construct
a finite cover by dominated charts satisfying the displayed upper, lower, and
mixed-Lipschitz bounds.  Critical, grazing, and zero-evidence observations form
separate closed graph-completed strata.  Prediction maps each chart into a
finite union of such charts and satisfies a common Lyapunov drift.
\end{theorem}

\begin{proof}
On a regular stratum, the observation map has a Jacobian minor bounded below.
The B1 exact-number coefficient and B2 actual-contact coarea formula give the
joint density \(k_\rho\) and its first source/state derivatives uniformly on a
compact chart.  Positivity of the reference density and compactness give
\(c_L,C_L\).  Integrating the derivative estimate in the observation variable
gives the mixed \(L^1(BL^*)\) Lipschitz bound.  Where the Jacobian rank changes,
stratify by the constant-rank theorem and use the corresponding Hausdorff
measure as the dominating measure.  Graph completion retains the incident
germs.  The B4 law semigroup and exponential moment estimate give the
prediction and Lyapunov assertions.
\end{proof}

\subsection{Stable disintegration}

Let \(\Pi(\rho,\cdot)\) be the law of the posterior random measure when
\(Y\) is sampled from \(g_\rho\lambda\).

\begin{theorem}[Integrated filter stability]
\label{thm:r16-c1-stability}
On every regular dominated chart,
\[
 \mathcal W_{d_{BL,W}\wedge1}
 \bigl(\Pi(\rho,\cdot),\Pi(\tilde\rho,\cdot)\bigr)
 \le C_L d_{\mathcal I}(\rho,\tilde\rho).
\]
More precisely, after an optimal coupling of observation values,
\[
 \int d_{BL,W}(B(\rho,y),B(\tilde\rho,y))
       (g_\rho\wedge g_{\tilde\rho})(y)d\lambda(y)
 \le {2\over c_L}
 \|k_\rho-k_{\tilde\rho}\|_{L^1_\lambda(BL_W^*)}.
\]
Thus posterior disintegration is continuous in the information topology,
not in the weak topology of joint laws.  Oscillatory graph sequences which
converge only weakly do not converge in \(d_{\mathcal I}\).
\end{theorem}

\begin{proof}
For positive finite measures \(\mu_y=k_\rho(\cdot,y)m\) and
\(\tilde\mu_y=k_{\tilde\rho}(\cdot,y)m\), write
\(g=\mu_y(1)\), \(\tilde g=\tilde\mu_y(1)\).  The elementary normalization
inequality
\[
 (g\wedge\tilde g)
 d_{BL,W}(\mu_y/g,\tilde\mu_y/\tilde g)
 \le2\|\mu_y-	ilde\mu_y\|_{BL,W}^*
\]
followed by integration proves the second estimate.  Couple the common part
of the evidence densities at the same \(y\) and the unmatched mass
arbitrarily; the evidence-marginal total variation is controlled by the same
mixed norm.  This gives the Wasserstein estimate.  In the oscillatory example,
the conditional density has order-one \(L^1\) oscillation, so the sequence is
not Cauchy in the declared topology.
\end{proof}

\subsection{Positive slices and singular strata}

Let \(T_\rho\) be the positive measure-current pair of a predicted law and let
\(O\) be a stratified smooth observation map.  The Federer slice of the
oriented current is used to calculate geometry; its mass measure is the
positive coefficient measure.

\begin{theorem}[Typed repeated observation]
\label{thm:r16-c1-slices}
For a finite sequence of constant-rank observation maps, iterated Federer
slices exist almost everywhere and agree with the slice by the product map.
At every stage the state records separately:
\begin{enumerate}[label=(\roman*)]
\item the positive unnormalized coefficient measure;
\item the oriented geometric current and its dimension;
\item the evidence density with respect to the stratum measure; and
\item the normalized posterior when the evidence is positive.
\end{enumerate}
Prediction and actual-contact insertion preserve this typed category.  No
positivity is asserted for an oriented current merely because its coefficient
measure is positive.
\end{theorem}

\begin{proof}
The slicing theorem for normal currents gives the oriented slice and its mass
bound on each constant-rank stratum.  Rokhlin disintegration of the positive
coefficient measure gives the positive conditional law.  Fubini proves
associativity for product observations.  B2's incoming flux is a positive
coefficient measure with a separately oriented contact graph; B4 prediction
is a positive push-forward.  Hence all four components remain well typed.
\end{proof}

\subsection{Finite determining coordinates chosen by tolerance}

For a compact information-state set \(K\) and \(\varepsilon>0\), compactness
provides a finite family
\(\Phi_\varepsilon=\{\phi_1^\varepsilon,\ldots,
\phi_{N(\varepsilon)}^\varepsilon\}\) in the weighted bounded-Lipschitz unit
ball such that equality within \(\varepsilon\) on those tests implies
\(d_{BL,W}\le3\varepsilon\) on \(K\).  Choose a measurable nearest-point
reconstruction from a finite \(\varepsilon\)-net.

\begin{theorem}[Tolerance-adapted finite belief games]
\label{thm:r16-c1-finite}
For every finite horizon and bounded Lipschitz reward there are tolerance-
selected finite determining families and reconstructions for which the
finite-coordinate controlled games converge uniformly on compact information
sets to the full belief game.  If \(V_n\) and \(V_n^\varepsilon\) are the
values after \(n\) observation steps, then
\[
 \sup_{\rho\in K}|V_n(\rho)-V_n^\varepsilon(\rho)|
 \le Cn\varepsilon+C\mathbb P\{\rho_j\notin K_L
 \text{ for some }j\le n\}.
\]
No property is required of the first \(N\) members of an arbitrarily ordered
dense sequence.
\end{theorem}

\begin{proof}
Compactness of \(K\) and the dual characterization of the bounded-Lipschitz
metric give the finite determining family.  Theorem~\ref{thm:r16-c1-stability}
propagates a one-step reconstruction error by at most \(C\varepsilon\).
Backward induction through the minimax Bellman recursion gives the first term.
The common Lyapunov drift controls exit from \(K_L\); take \(L\to\infty\) and
then \(\varepsilon\downarrow0\).
\end{proof}

\subsection{Dynamic programming and regular posterior asymptotics}

\begin{theorem}[Belief-state dynamic programming]
\label{thm:r16-c1-dpp}
On the union of dominated and graph-completed singular charts, prediction
followed by the typed observation kernel constructs the exact controlled law.
For compact relaxed controls and bounded Lipschitz rewards the full belief
game satisfies dynamic programming.  One-time preparation, reward-only
feedback, and adaptive actual-contact likelihood control retain their distinct
timing and Hamiltonians.
\end{theorem}

\begin{proof}
The positive joint kernel and its Rokhlin disintegration define a measurable
transition on every chart; the chart label is included in the state.  The
Ionescu--Tulcea construction gives the controlled process.  On regular charts
Theorem~\ref{thm:r16-c1-stability} is Feller; graph-completed singular strata
are reached by their explicit kernels and are treated by lower-semicontinuous
relaxation.  Compactness of relaxed controls and measurable selection give the
Bellman recursion.  B2/B4 identify the actual-contact Hamiltonian; the static
preparation multiplier is not reoptimized inside the transition game.
\end{proof}

Let \(\theta\) parametrize a compact regular dominated family.  B1 and B3 now
supply exact finite centering, uniform third pressure derivatives, positive
information, and process Gaussianity.

\begin{theorem}[Controlled LAN and Bernstein--von Mises]
\label{thm:r16-c1-bvm}
Uniformly over compact predictable strategies and bounded local alternatives,
\[
 \log {dP_{\theta+h/\sqrt\mu}\over dP_\theta}
 =h^T\Delta_\mu-{1\over2}h^TI(\theta)h+o_{P_\theta}(1),
 \qquad
 \Delta_\mu\Rightarrow N(0,I(\theta)).
\]
For a prior with positive continuous density at \(\theta\), the posterior
centered at an efficient estimator converges in total variation on compact
sets to \(N(0,I(\theta)^{-1})\).  The result is asserted only on the regular
dominated chart proved in Theorem~\ref{thm:r16-c1-chart}; singular strata have
their separately typed posterior law.
\end{theorem}

\begin{proof}
Taylor-expand the exact finite-volume log normalizer at
\(\theta+h/\sqrt\mu\).  B1/B3 give uniform first, second, and third derivative
bounds and the Gaussian score process, proving LAN.  Identifiability gives
uniform exponentially consistent tests through the pressure rate gap.  The
inserted local coefficient supplies the likelihood lower bound on the regular
chart.  Standard Laplace expansion, with the exact finite centering retained,
gives total-variation Bernstein--von Mises convergence.
\end{proof}

\begin{theorem}[Round-sixteen C1 closure]
\label{thm:r16-c1-main}
Belief dynamics is Feller in a model-derived dominated information topology,
not under weak convergence of joint laws.  Positive coefficient measures,
oriented slices, evidence, and posteriors are separately typed; finite
coordinates are chosen by tolerance; singular and zero-evidence strata remain
explicit; and the regular controlled statistical limit uses the now-proved
B1--B3 inputs.
\end{theorem}

\begin{proof}
Combine Theorems~\ref{thm:r16-c1-chart},
\ref{thm:r16-c1-stability}, \ref{thm:r16-c1-slices},
\ref{thm:r16-c1-finite}, \ref{thm:r16-c1-dpp}, and
\ref{thm:r16-c1-bvm}.
\end{proof}
